\documentclass[journal]{IEEEtran}

\usepackage{hyperref}
\usepackage{amsmath, amssymb, amsthm}
\usepackage[lined,boxed,ruled,linesnumbered]{algorithm2e}
\usepackage{graphicx}
\usepackage{subcaption}
\usepackage{booktabs}
\usepackage{multirow}
\usepackage{tabularx}
\usepackage[dvipsnames]{xcolor}
\usepackage{cleveref}
\usepackage{enumitem}
\usepackage{cite}
\usepackage{adjustbox}
\usepackage{array}
\usepackage{makecell}
\usepackage{pifont}

\newtheorem{theorem}{Theorem}
\newtheorem{proposition}{Proposition}
\newtheorem{definition}{Definition}
\newtheorem{remark}{Remark}
\newcommand{\cmark}{\ding{51}}
\newcommand{\xmark}{\ding{55}}

\begin{document}

\title{Parameter-Efficient Quantum Graph Neural Networks for Graph Classification via QFI-Guided Geometric Pruning}

\author{%
  \IEEEauthorblockN{Cat Khanh Tran,~\IEEEmembership{}
    Thai Quang Nguyen Pham,~\IEEEmembership{}
    Van-Linh Nguyen,~\IEEEmembership{Member,~IEEE}}
  \thanks{C.-K. Tran, T.-Q.~N. Pham, and V.-L. Nguyen are with the Department
    of Computer Science and Information Engineering, National Chung Cheng
    University, Chiayi, Taiwan (e-mail: 614410145@office365.ccu.edu.tw , nvlinh@cs.ccu.edu.tw).}
  \thanks{Corresponding author: Van-Linh Nguyen.}
  \thanks{Source code: \url{https://github.com/jeremy0612/Q-Drop-Integration}}
}

\markboth{IEEE TRANSACTIONS ON QUANTUM ENGINEERING}%
{Tran \MakeLowercase{\textit{et al.}}: QFI-Pruned Quantum Graph Neural Networks}

\maketitle

%% ─────────────────────────────────────────────────────────────────────────────
\begin{abstract}
Graph classification is a fundamental graph-level learning problem that assigns a label to an entire graph rather than to individual nodes or edges. It remains challenging because graphs are irregular, non-Euclidean structures whose node semantics and global topology must be captured efficiently. Classical graph neural networks (GNNs) achieve strong accuracy but often require large parameter budgets that limit deployment on resource-constrained platforms. To address this, we propose two parameter-efficient quantum graph neural networks—the Quantum Graph Convolutional Network (QGCN) and Quantum Graph Attention Network (QGAT)—which use variational quantum circuit (VQC) encoders to produce compact quantum node representations for graph-level prediction. However, quantum neural networks remain vulnerable to training instability caused by barren plateaus and gate co-adaptation. We therefore introduce QFI-Drop, a geometry-aware pruning algorithm that scores each variational gate by its spectral leverage on the data-averaged Quantum Fisher Information Matrix (QFIM): gates whose energy lies predominantly in the null eigendirections of the QFIM are frozen, and the surviving gates are updated by a unified quantum-natural-gradient step preconditioned on the QFIM sub-manifold. Unlike heuristic gate-dropout schemes, QFI-Drop uses the intrinsic information geometry of the circuit to identify which parameters actually contribute to the trainable landscape, mitigating barren plateaus and co-adaptation at their source. To preserve permutation equivariance, the QFIM probe and the resulting freeze mask are shared across all node-level circuit evaluations within a mini-batch, so pruning depends only on graph content and not on node ordering. Evaluated on five benchmarks—MUTAG, PROTEINS, NCI1, IMDB-Binary, and IMDB-Multi. The framework achieves competitive accuracy against classical baselines while using only 3,411–9,031 parameters in total, corresponding to more than 97\% parameter reduction.


\end{abstract}

\begin{IEEEkeywords}
Quantum graph neural network, quantum Fisher information, geometric pruning,
quantum natural gradient, graph classification, variational quantum circuit,
quantum machine learning, barren plateau, parameter efficiency
\end{IEEEkeywords}

\maketitle

%% ─────────────────────────────────────────────────────────────────────────────
\section{Introduction}
\label{sec:intro}

Graph classification---assigning a categorical label to an entire graph---is
a fundamental task in machine learning with high-value applications spanning
computational chemistry~\cite{MUTAG1991}, structural
biology~\cite{PROTEINS2003}, social network
analysis~\cite{IMDB2015}, and blockchain security~\cite{Zhuang2020IJCAI,
Liu2021TKDE,Qian2023WWW}. The problem requires models to capture both
local node-level semantics and global graph topology, making it
substantially harder than node-level or edge-level tasks.

Classical graph neural networks (GNNs) have become the dominant paradigm
for graph classification. GCN~\cite{Kipf2017} demonstrated that stacked
spectral convolution layers learn powerful graph representations.
GAT~\cite{Velickovic2018} and GATv2~\cite{Brody2022ICLR} introduced
attention-weighted aggregation, improving differentiation of neighbor
relevance. GIN~\cite{Xu2019GIN} established theoretical limits by matching
the 1-Weisfeiler--Leman (WL) expressiveness test. Higher-order approaches,
including OSAN~\cite{Bause2023ICML}, CIN~\cite{Bodnar2021CIN},
DropGNN~\cite{Papp2021NeurIPS}, and GRATIN~\cite{Abbahaddou2025ICML},
push beyond 1-WL at the cost of significantly larger parameter budgets.
Even lightweight configurations typically carry 50K--500K parameters,
which is prohibitive for resource-constrained deployment scenarios such
as edge computing nodes, embedded biosensors, blockchain validator
gateways, and on-device medical diagnostics.

Quantum machine learning (QML)~\cite{Cerezo2021,
Schuld2021} opens a fundamentally different design space. Parameterized
quantum circuits (PQCs)~\cite{Benedetti2019} exploit superposition and
entanglement to implement rich nonlinear transformations in an exponentially
large Hilbert space while requiring far fewer classically trainable
parameters than neural networks of comparable expressive capacity.
Early graph-level QML work includes equivariant quantum graph
circuits~\cite{Mernyei2022ICML}, the quantum neural tangent kernel for
graphs~\cite{Tang2022NeurIPS}, and quantum walk kernel methods such as
QBMK~\cite{Bai2024ICML}. These pioneering works established that
quantum-inspired approaches are viable for graph-structured data.

Building on this foundation, our group introduced two unpublished hybrid
architectures specifically designed for quantum graph classification:
the Quantum Graph Convolutional Network (QGCN) and the
Quantum Graph Attention Network (QGAT). QGCN applies a
node-wise variational quantum circuit (VQC) to produce compact quantum
node embeddings and propagates them via classical GCN aggregation,
achieving competitive smart contract vulnerability detection accuracy
with only $\approx$2,032 parameters---a 98\% reduction relative to
classical GNN baselines~\cite{Zhuang2020IJCAI,Liu2021IJCAI,Liu2021TKDE,
Qian2023WWW} that carry 66K--263K parameters. QGAT further replaces
classical aggregation with a shared quantum attention circuit for edge
scoring, operating with only 1,828 parameters while achieving
100\% F1 on integer overflow/underflow
detection. These results demonstrate the parameter
efficiency of quantum graph encoders, but they also reveal a
fundamental gap: \emph{no systematic regularization exists for
graph-structured quantum encoders}.

This gap is critical. As quantum circuit width or depth grows,
gradient magnitudes vanish exponentially---the barren plateau
phenomenon~\cite{McClean2018}---stalling optimization. Even in the
shallow circuits used by QGCN and QGAT, unconstrained training leads
to \emph{gate co-adaptation}: a small subset of high-gradient gates
consistently dominates the representation, suppressing the contribution
of remaining parameters and reducing generalization. In classical deep
learning, dropout~\cite{Srivastava2014} and magnitude-based
pruning~\cite{Han2015,LeCun1990OBD} are well-established remedies to
analogous co-adaptation and redundancy problems. Analogous strategies
for quantum circuits remain largely unexplored, with limited work such
as Entangle Dropout~\cite{EntangleDropout} and layerwise
learning~\cite{Skolik2021} addressing only narrow aspects.

Our group's prior published work Q-Drop~\cite{QDrop2025ICC} took an
important step by introducing two complementary quantum regularization
techniques for quantum orthogonal neural networks on \emph{image}
classification tasks: (1) \emph{dynamic circuit dropout} (circuit
flip-flop), which stochastically deactivates Reconfigurable Beam
Splitter (RBS) gate subsets via a Bernoulli mask at each training
epoch; and (2) \emph{statistical pruning} (SP), which accumulates
per-gate gradient magnitudes over a sliding window and selectively
freezes low-impact gates based on the resulting importance distribution.
Q-Drop achieved 94.3\% accuracy on Pneumonia-MNIST and 98.7\% on
MNIST, outperforming qOrthoNN~\cite{Kerenidis2020QONN}, qCNN, and
the parameter-shift rule baseline. However, Q-Drop was designed
for flat (non-graph) inputs: the VQC processes a single sample at a
time, and the dropout mask is applied globally.


In this paper, we close these gaps by proposing \textbf{Geometric Adapted QFI-pruning on QGCN and QGAT}, the first quantum graph neural networks with systematic
regularization, submitted to IEEE Transactions on Quantum Engineering.
The central technical innovations are: (1) \emph{graph-consistent
masking}, which samples a single gate mask per mini-batch and applies
it uniformly to all nodes, preserving permutation equivariance; and
(2) \emph{graph-adapted statistical pruning}, which aggregates gradient
magnitudes across all node-level circuit evaluations within a sliding
window, enabling robust gate importance estimation under graph-size
variance. We evaluate on five standard graph classification benchmarks
and demonstrate competitive accuracy against classical GNNs carrying
50K--500K parameters, while using only 3,411--9,031 parameters total.
The contributions of this paper are as follows:
\begin{itemize}
    \item \textbf{Two novel QGNN architectures:} QGCN and
    QGAT are proposed for the first time on general graph classification
    benchmarks (MUTAG, PROTEINS, NCI1, IMDB-Binary, IMDB-Multi).

    \item \textbf{Graph-adapted geometric pruning:} Each gate is scored by its
    spectral leverage on the data-averaged Quantum Fisher Information Matrix,
    probed over all nodes in a mini-batch. Null-eigendirection gates are
    frozen and survivors updated by a quantum-natural-gradient step on the
    QFIM sub-manifold—mitigating barren plateaus at their geometric source
    while the shared per-batch mask preserves permutation equivariance.

    \item \textbf{Comprehensive benchmark and ablation:} We compare
    against DropGNN, GATv2, OSAN, GRDL, SE2P, GRATIN, and QBMK,
    demonstrating $>$97\% parameter reduction at competitive accuracy;
    ablation confirms Quantum Network with suitable methods are individually necessary and mutually complementary.
    \href{https://github.com/jeremy0612/Q-Drop-Integration}{\textcolor{blue}{Q-Drop}}.
    \footnote{\url{https://github.com/jeremy0612/Q-Drop-Integration}}
\end{itemize}

The remainder of this paper is organized as follows.
Section~\ref{sec:related} reviews related work.
Section~\ref{sec:background} provides background on graph classification
and quantum circuit primitives.
Section~\ref{sec:arch} presents the QGCN and QGAT architectures in detail.
Section~\ref{sec:method} introduces the Q-Drop integration framework.
Section~\ref{sec:eval} reports experimental results and analysis.
Section~\ref{sec:conclusion} concludes the paper.

%% ─────────────────────────────────────────────────────────────────────────────
\section{Related Work}
\label{sec:related}

\subsection{GNN-Based Graph Classification}

Graph classification has been studied extensively since the introduction of
graph kernels~\cite{Vishwanathan2010,Shervashidze2011WL} and the subsequent
wave of deep learning methods. GCN~\cite{Kipf2017} framed spectral convolution
as localized first-order approximations, enabling efficient feature propagation.
GraphSAGE~\cite{Hamilton2017} introduced inductive sampling-based aggregation,
scaling to large unseen graphs. GIN~\cite{Xu2019GIN} provably maximized
discriminative power by using injective aggregation, matching the 1-WL graph
isomorphism test. GAT~\cite{Velickovic2018} brought attention-weighted
neighborhood aggregation, and GATv2~\cite{Brody2022ICLR} corrected GAT's
static attention limitation by computing dynamic scores through a
pre-nonlinearity transformation, strictly increasing expressiveness.

Beyond 1-WL expressiveness, higher-order approaches have been pursued
extensively. OSAN~\cite{Bause2023ICML} introduced ordered subgraph aggregation
for outerplanar graphs, achieving 93.6\% on MUTAG. CIN~\cite{Bodnar2021CIN}
used cellular sheaf structures, GSN~\cite{Bouritsas2022} leveraged subgraph
isomorphism counting, and PathNN~\cite{Michel2023PathNN} integrated path-level
information, each exceeding 1-WL at the cost of significantly increased
computational and memory overhead. Stochastic depth approaches such as
DropGNN~\cite{Papp2021NeurIPS} execute multiple GNN runs over graphs with
randomly dropped nodes, improving expressiveness by inducing ensemble effects
over diverse structural views of each graph.

Recent augmentation-based approaches target distribution shift and
out-of-distribution generalization. GRDL (NeurIPS 2024) combines
graph-level regularization with differential learning. SE2P~\cite{Saber2024CIKM}
introduces scalable graph perturbation in the preprocessed feature space.
GRATIN~\cite{Abbahaddou2025ICML} fits Gaussian Mixture Models to entire
training-set graph representations and samples augmented views in latent space,
achieving strong OOD generalization. ANGC (AAAI 2026) proposes an anchor-driven
Nystr\"om framework for deep graph-level clustering. Despite their accuracy
these methods all operate entirely in the classical domain and typically
carry 50K--500K+ parameters.

\subsection{Quantum Graph Learning}

Quantum approaches to graph-structured learning have attracted growing
research interest as quantum hardware matures. QBMK~\cite{Bai2024ICML}
uses continuous-time quantum walk (CTQW) Shannon entropy kernels to compute
graph similarity, capturing global and local structural statistics via quantum
walk statistics on aligned vertices. While QBMK provides an interesting
quantum perspective on graph similarity, it lacks neural network feature
learning, employs no attention mechanism, uses no dropout regularization,
and relies on a classical SVM readout, limiting its end-to-end trainability
and expressiveness.

Equivariant quantum graph circuits~\cite{Mernyei2022ICML} provide theoretical
foundations for symmetry-preserving quantum graph operations by constructing
circuits that respect graph automorphisms. GraphQNTK~\cite{Tang2022NeurIPS}
studied the quantum neural tangent kernel for graph data, providing theoretical
insights into the training dynamics of quantum graph models in the
infinite-width limit. However, none of these methods have been evaluated at
scale on standard graph classification benchmarks, and none incorporate
systematic regularization strategies.

\subsection{Quantum Circuit Regularization}

Regularization for parameterized quantum circuits is a nascent but rapidly
developing field. The barren plateau problem~\cite{McClean2018} describes
the exponential vanishing of gradients in deep or wide quantum circuits,
motivating the need for structural and training-time interventions.
Entangle Dropout~\cite{EntangleDropout} applied classical dropout ideas
directly to entanglement gates, showing that randomly removing entangling
operations during training can provide modest robustness improvements on
small circuits. Layerwise learning~\cite{Skolik2021} proposed training
quantum circuits one layer at a time to avoid barren plateaus during
initialization. Noise-induced barren plateaus~\cite{Wang2021Noise} studied
how hardware noise exacerbates gradient vanishing, motivating noise-resilient
training procedures.

Our Q-Drop~\cite{QDrop2025ICC} provides the most comprehensive regularization
framework for quantum neural networks to date, combining two orthogonal
mechanisms: dynamic dropout via a circuit flip-flop that stochastically
deactivates gate subsets, and statistical pruning that identifies low-impact
gates by tracking accumulated gradient magnitudes. Q-Drop achieved 94.3\%
accuracy on Pneumonia-MNIST and 98.7\% on standard MNIST, outperforming
qOrthoNN, the Parameter Shift Rule baseline, and qCNN. Crucially, Q-Drop
has never been applied to graph-structured data, and the graph-consistency
requirement---that all nodes within a graph must experience a coherent
circuit structure per training step---represents a non-trivial design
challenge that we address in this paper.

\subsection{Summary and Positioning}

Table~\ref{tab:related_comparison} summarizes the key properties of the
most closely related methods. Our proposed framework is the only approach
that simultaneously provides: hybrid classical-quantum encoding, graph
attention, dropout regularization, gradient-based pruning, and evaluation
on standard multi-dataset graph benchmarks.

{\renewcommand{\arraystretch}{1.2}
\begin{table*}[t]
  \centering
  \caption{Comparison of closely related methods across key design dimensions.
    \cmark~= present, \xmark~= absent, P = partial.}
  \label{tab:related_comparison}
  \resizebox{\linewidth}{!}{%
  \begin{tabular}{l l l c c c c c c}
    \toprule
    Method & Venue & Year &
    \makecell{Neural\\Learning} &
    \makecell{Graph\\Attention} &
    \makecell{Quantum\\Encoder} &
    \makecell{Dropout\\Reg.} &
    \makecell{Gradient\\Pruning} &
    \makecell{Multi-dataset\\Benchmark} \\
    \midrule
    DropGNN~\cite{Papp2021NeurIPS}   & NeurIPS & 2021 & \cmark & \xmark & \xmark & P & \xmark & \cmark \\
    GATv2~\cite{Brody2022ICLR}       & ICLR    & 2022 & \cmark & \cmark & \xmark & \xmark & \xmark & \cmark \\
    QBMK~\cite{Bai2024ICML}          & ICML    & 2024 & \xmark & \xmark & P      & \xmark & \xmark & P \\
    GRATIN~\cite{Abbahaddou2025ICML} & ICML    & 2025 & \cmark & \xmark & \xmark & \xmark & \xmark & \cmark \\
    \textbf{Ours}          & --      & 2026 & \cmark & \cmark & \cmark & \cmark & \cmark & \cmark \\
    \bottomrule
  \end{tabular}}
\end{table*}}

%% ─────────────────────────────────────────────────────────────────────────────
\section{Background}
\label{sec:background}

\subsection{Graph Classification Problem}

Let $G = (V, E, \mathbf{X})$ denote a graph with node set $V$, edge set $E$,
and node feature matrix $\mathbf{X} \in \mathbb{R}^{|V| \times d}$, where
$d$ is the feature dimension. In graph classification, we are given a labeled
training set $\mathcal{D} = \{(G_i, y_i)\}_{i=1}^N$ and seek to learn a
parametric predictor $f_\Theta: G \to \mathcal{Y}$ that maps each graph to
a class label $y_i \in \mathcal{Y}$. The standard training objective is to
minimize the empirical risk:
\begin{equation}
  \min_\Theta \;\frac{1}{N}\sum_{i=1}^N \mathcal{L}\!\left(f_\Theta(G_i), y_i\right),
  \label{eq:obj}
\end{equation}
where $\mathcal{L}$ is the cross-entropy loss for multi-class settings or
binary cross-entropy for binary classification.

A typical GNN-based $f_\Theta$ consists of $L$ message-passing layers that
iteratively update node representations, followed by a global pooling
operator and a classification head. The representation of node $v$ after
layer $\ell$ is:
\begin{equation}
  \mathbf{h}_v^{(\ell)} = \phi\!\left(\mathbf{h}_v^{(\ell-1)},\;
  \text{AGG}\!\left\{\mathbf{h}_u^{(\ell-1)} : u \in \mathcal{N}(v)\right\}\right),
  \label{eq:gnn}
\end{equation}
where $\phi$ is a learnable update function, $\text{AGG}$ is a neighborhood
aggregation operator, and $\mathcal{N}(v)$ denotes the neighbors of $v$.
After $L$ layers, a readout function $R$ produces the graph embedding:
$\mathbf{g} = R\!\left(\{\mathbf{h}_v^{(L)}\}_{v \in V}\right)$, typically
global mean or sum pooling.

\subsection{Variational Quantum Circuits}

A variational quantum circuit (VQC) is a parameterized quantum program
consisting of alternating layers of fixed entangling gates and trainable
single-qubit rotations~\cite{Cerezo2021}. Given an $n_q$-qubit system,
the circuit takes an input state $|0\rangle^{\otimes n_q}$ and produces
an output state $|\psi(\theta)\rangle = U(\theta)|0\rangle^{\otimes n_q}$,
where $U(\theta)$ is the product of all gate operations parameterized by
$\theta \in \mathbb{R}^P$.

\emph{Angle embedding} encodes a classical input vector $\mathbf{z} \in
\mathbb{R}^{n_q}$ into the quantum state by applying $R_Y(z_i)$ rotations:
$|\psi(\mathbf{z})\rangle = \bigotimes_{i=1}^{n_q} R_Y(z_i)|0\rangle$,
where $R_Y(\alpha) = e^{-i\alpha Y/2}$ is a rotation about the $Y$-axis.
The variational ansatz then applies trainable $R_Y(\theta_i)$ and
$R_Z(\phi_i)$ rotations on each qubit, interspersed with a ring
CNOT entanglement pattern. Pauli-Z measurements extract classical
outputs:
$g_i = \langle\psi(\theta)|Z_i|\psi(\theta)\rangle \in [-1, 1]$.

\emph{Parameter-shift rule}~\cite{Mitarai2018} enables exact gradient
computation for VQCs on quantum hardware:
\begin{equation}
  \frac{\partial \mathcal{L}}{\partial \theta_i}
  = \frac{1}{2}\left[\mathcal{L}(\theta_i + \tfrac{\pi}{2})
  - \mathcal{L}(\theta_i - \tfrac{\pi}{2})\right],
  \label{eq:psr}
\end{equation}
making VQCs fully differentiable and compatible with classical
backpropagation frameworks.

%% ─────────────────────────────────────────────────────────────────────────────
\section{QGCN and QGAT: Base Architectures}
\label{sec:arch}

\subsection{Overview}

Both QGCN and QGAT follow the same high-level hybrid design: a classical
graph preprocessing stage normalizes the input graph, a stacked quantum
encoder transforms node features into compact quantum embeddings, and a
lightweight classification head maps the pooled graph embedding to a
class label. The key distinction is in the aggregation strategy: QGCN
uses classical spectral aggregation derived from GCN, while QGAT replaces
this with a learnable quantum attention mechanism for edge scoring.

\subsection{QGCN: Quantum Graph Convolutional Network}
\label{ssec:qgcn}

\subsubsection{Graph Normalization}

Following, we apply a merge-and-redirect normalization
procedure that eliminates secondary and fallback nodes by merging their
features into the nearest major nodes by graph distance and redirecting
incident edges accordingly. This reduces structural variance across graphs
of differing styles while preserving the execution-relevant topology,
producing a normalized graph $\bar{G} = (\bar{V}, \bar{E})$ with enriched
major node features.

\subsubsection{Node-wise Quantum Feature Extraction}

At layer $\ell$, each node $v \in \bar{V}$ has a representation
$\mathbf{h}_v^{(\ell)} \in \mathbb{R}^{d_\ell}$. A learnable
projection matrix $W_r^{(\ell)} \in \mathbb{R}^{d_\ell \times n_q}$
maps this to an $n_q$-dimensional quantum embedding:
\begin{equation}
  \mathbf{z}_v^{(\ell)} = W_r^{(\ell)}\,\mathbf{h}_v^{(\ell)}
  \in \mathbb{R}^{n_q}.
  \label{eq:proj}
\end{equation}
The VQC encodes $\mathbf{z}_v^{(\ell)}$ via angle embedding and applies
a shallow variational ansatz with trainable single-qubit rotations
$\{R_Y(\theta_i), R_Z(\phi_i)\}$ and a ring CNOT entanglement pattern.
Pauli-Z measurements produce the quantum embedding:
\begin{equation}
  \mathbf{g}_v^{(\ell)} = Q_\phi\!\left(\mathbf{z}_v^{(\ell)}\right)
  = \left[\langle Z_0\rangle, \ldots, \langle Z_{n_q-1}\rangle\right]
  \in \mathbb{R}^{n_q}.
  \label{eq:vqc_qgcn}
\end{equation}
The same VQC template $Q_\phi$ with shared parameters $\phi$ is applied
independently to every node in the graph, keeping the parameter count
independent of graph size.

\subsubsection{Classical Neighborhood Aggregation}

The quantum embeddings are aggregated via the classical GCN rule:
\begin{equation}
  \mathbf{H}^{(\ell+1)} = \sigma\!\left(\hat{A}\,\mathbf{G}^{(\ell)}\right),
  \label{eq:gcn_agg}
\end{equation}
where $\mathbf{G}^{(\ell)}$ stacks all $\mathbf{g}_v^{(\ell)}$ row-wise,
$\hat{A} = \tilde{D}^{-1/2}\tilde{A}\tilde{D}^{-1/2}$ is the symmetrically
normalized adjacency matrix with $\tilde{A} = A + I$, $\tilde{D}$ is the
degree matrix of $\tilde{A}$, and $\sigma$ is a non-linearity (LeakyReLU).

\subsubsection{Classification Heads}

After $L$ layers, global mean pooling produces the graph embedding:
$\mathbf{g} = \frac{1}{|\bar{V}|}\sum_{v \in \bar{V}}\mathbf{h}_v^{(L)}$.
We evaluate three classifier head variants on top of $\mathbf{g}$:

\textit{Linear classifier:} A simple affine transformation
$\hat{y} = \sigma(\mathbf{w}^\top\mathbf{g} + b)$, where $\sigma$ is
sigmoid for binary tasks. This head is intentionally lightweight,
ensuring that performance improvements can be attributed to the quantum
encoder rather than classifier complexity.

\textit{Matrix Product State (MPS) classifier:} Adopted from~\cite{Bhatia2019MPS},
this head encodes $\mathbf{g}$ into $n_q$ qubits via angle embedding and
applies a chain-structured entangling circuit that couples adjacent qubits.
The readout is obtained from a Pauli-Z expectation on a designated output
qubit.

\textit{Tree Tensor Network (TTN) classifier:} Adopted
from~\cite{Grant2018Hierarchical}, this head uses a binary tree circuit
where qubits are progressively merged via pairwise entangling operations.
The hierarchical structure naturally captures multi-scale correlations in
the graph embedding.

In all experiments we report results for the linear classifier as the
primary variant, which consistently achieves the best performance,
suggesting that the quantum encoder already produces linearly separable
representations for the tasks considered.
The node-wise quantum encoder is illustrated in
Fig.~\ref{fig:node-vqc}. This circuit is shared by both QGCN and QGAT:
each node is processed independently by the same VQC template, while
the trainable quantum parameters are shared across all node evaluations.
This design makes the quantum parameter count independent of the number
of nodes in the graph.

\begin{figure}[!t]
    \centering
    \includegraphics[width=\columnwidth]{QGCN.pdf}
    \caption{
    Node-wise variational quantum circuit encoder shared by QGCN and QGAT.
    Each node representation $h_v^{(\ell)}$ is first projected by
    $W_r^{(\ell)}$ into an $n_q$-dimensional vector $z_v^{(\ell)}$.
    The vector is encoded through angle embedding using $R_Y(z_i)$ rotations,
    followed by trainable $R_Y(\theta_i)$ and $R_Z(\phi_i)$ rotations and a
    CNOT entangling layer. Pauli-$Z$ measurements produce the quantum node
    embedding $g_v^{(\ell)} = [\langle Z_0\rangle,\ldots,\langle Z_{n_q-1}\rangle]$.
    The same circuit template and parameters are shared across all nodes,
    keeping the quantum parameter count independent of graph size.
    }
    \label{fig:node-vqc}
\end{figure}
\subsubsection{Node-wise Quantum Feature Extraction}

QGAT uses the same node-wise VQC as QGCN to produce quantum node
embeddings $\mathbf{h}_v^{(0)} \in \mathbb{R}^{n_q}$ at the first layer,
and updates them to $\mathbf{h}_v^{(\ell)} \in \mathbb{R}^{n_q}$ across
subsequent layers. The circuit architecture is identical to
Eq.~\eqref{eq:vqc_qgcn} with the same angle embedding and variational
ansatz design.

\subsubsection{Quantum Attention Mechanism}

For each directed edge $(u, v) \in \bar{E}$, QGAT computes a scalar
attention logit via a quantum attention circuit. An edge-specific input
vector is formed by concatenation and low-dimensional projection:
\begin{equation}
  \mathbf{z}_{uv} = \mathbf{W}_a\,[\mathbf{h}_u \| \mathbf{h}_v]
  \in \mathbb{R}^{n_q},
  \label{eq:edge_proj}
\end{equation}
where $\mathbf{W}_a \in \mathbb{R}^{n_q \times 2n_q}$ is a trainable
projection. The attention circuit (illustrated in Fig.~\ref{fig:qgat_circuit})
injects $\mathbf{z}_{uv}$ as data-dependent rotations $R_Y(z_i)$ on qubit
$i$, then applies shared trainable rotations $\{R_Y(\theta_i), R_Z(\phi_i)\}$
on each qubit and a readout rotation $R_Y(\omega)$ on the last qubit.
The scalar attention logit is obtained from the Pauli-Z expectation on
the last qubit:
\begin{equation}
  o_{uv} = \mathcal{Q}(\mathbf{z}_{uv}; \Theta) = \langle Z_{n_q-1}\rangle
  \in [-1, 1].
  \label{eq:qgat_logit}
\end{equation}
The trainable parameters $\Theta = \{\theta_i, \phi_i, \omega\}$ are
\textit{shared across all edges}, so the parameter count scales only with
$n_q$ and not with graph size.

\subsubsection{Attention-Weighted Aggregation}

The attention weights are computed by applying softmax over all incoming
edges to each node:
\begin{equation}
  \alpha_{uv} = \frac{\exp(o_{uv} / \tau)}{\sum_{u' \in \mathcal{N}(v)}
  \exp(o_{u'v} / \tau)},
  \label{eq:softmax_att}
\end{equation}
where $\tau$ is a temperature hyperparameter. Node embeddings are then
updated by attention-weighted aggregation:
\begin{equation}
  \mathbf{h}_v^{(\ell+1)} = \sigma\!\left(\sum_{u \in \mathcal{N}(v)}
  \alpha_{uv}\,\mathbf{h}_u^{(\ell)}\right).
  \label{eq:qgat_agg}
\end{equation}
This aggregation remains fully classical for scalability, while the
quantum circuit handles the attention scoring. After $L$ layers, global
mean pooling and a linear classifier produce $\hat{y}$ as in QGCN.

\begin{figure}[t]
  \centering
  % Placeholder: replace with actual tikz or pdf figure
  \includegraphics[width=\linewidth]{QGAT.pdf}
  \caption{Quantum attention circuit for QGAT ($n_q = 8$ qubits). Data-dependent
    rotations $R_Y(z_i)$ encode the edge feature vector. Trainable rotations
    $R_Y(\theta_i)$, $R_Z(\phi_i)$, and a readout rotation $R_Y(\omega)$ on
    the last qubit produce a scalar attention logit $\langle Z_7 \rangle$.
    All parameters $\{\theta_i, \phi_i, \omega\}$ are shared across edges.}
  \label{fig:qgat_circuit}
\end{figure}

%% ─────────────────────────────────────────────────────────────────────────────
\section{Proposed QFI-Pruning Framework for Graph Networks}
\label{sec:method}

\begin{figure*}[t]
  \centering
  \includegraphics[width=1\linewidth]{images/qfi_vs_sp_circuit.drawio.pdf}
  \caption{Original statistical pruning (a) vs.\ the proposed
    graph-adapted QFI-pruning (b) on the shared QGCN/QGAT node-embedding
    circuit (AngleEmbedding, variational $R_Y$ layer, CNOT ring).
    (a)~SP scores each gate by its gradient magnitude accumulated over a
    sliding window of $m$ steps, freezes gates below the pruning-ratio
    threshold $\rho$, and applies a plain gradient step to the survivors.
    (b)~QFI-pruning probes the data-averaged QFIM once per epoch,
    eigendecomposes it to isolate the informative eigenspace
    ($\lambda_i \geq \kappa\lambda_{\max}$), scores each gate by its
    spectral leverage $e_j$, freezes gates with $e_j < \tau$, and
    preconditions the survivors with the quantum-natural-gradient step of
    Eq.~\eqref{eq:sp_anneal}. The two criteria can disagree: a gate with a
    large but flat gradient---its energy confined to null
    eigendirections---is kept by SP yet frozen by QFI-pruning, whereas a
    small-gradient gate with high leverage is frozen by SP yet kept and
    QNG-rescaled. In both variants a single mask is shared by every
    node-level evaluation in the mini-batch, preserving permutation
    equivariance; at inference all gates are active.}
  \label{fig:framework}
\end{figure*}

This section presents the two core technical contributions of this work:
graph-consistent QFIM probing and graph-adapted QFI-pruning.
Fig.~\ref{fig:framework} contrasts the proposed criterion with the
original statistical pruning of~\cite{QDrop2025ICC}; the training
pipeline is summarized in Algorithm~\ref{alg:qdrop_graph}.

\subsection{Motivation: Training Instabilities in QGNNs}

Before presenting the technical details, we briefly analyze the sources
of training instability in quantum graph networks. When the same VQC
parameters $\phi$ are applied to all nodes independently, the gradients
with respect to $\phi$ are computed as:
\begin{equation}
  \nabla_\phi \mathcal{L} = \frac{1}{|B|}
  \sum_{G \in B}\sum_{v \in V_G} \frac{\partial \mathcal{L}_G}{\partial g_{v,i}}
  \cdot \frac{\partial g_{v,i}}{\partial \phi_i},
  \label{eq:total_grad}
\end{equation}
where $B$ is the current mini-batch and $V_G$ is the node set of graph $G$.
Two pathological behaviors emerge:

\textit{Gate co-adaptation:} When certain gate parameters consistently
dominate the gradient signal across many graph instances, other gates
become redundant and receive negligible updates. The model effectively
learns to use only a subset of its quantum degrees of freedom.

\textit{High gradient variance:} Because $|V_G|$ varies significantly
across graphs (e.g., from 17 nodes in MUTAG to over 600 in PROTEINS),
the magnitude of $\nabla_\phi \mathcal{L}$ fluctuates wildly between
mini-batches, destabilizing learning rate selection and convergence.

Both pathologies are geometric: they manifest as a QFIM whose spectrum
collapses onto a few dominant eigendirections while the remainder become
null. QFI-pruning addresses them jointly---freezing gates that live in the
null space and steering the survivors along the informative sub-manifold
via a quantum-natural-gradient step.

\subsection{Graph-Consistent QFIM Probing}
\label{ssec:qdrop}
The original Q-Drop~\cite{QDrop2025ICC} operates on flat inputs where a
single circuit processes one sample at a time. Since this method did not bring major effects on graph issues, we came up with new geometric-oriented improvement for the algorithm.
QFI-pruning is driven by the Quantum Fisher Information Matrix (QFIM),
the Fubini--Study metric of the node-embedding circuit, which measures
how distinguishable the output state becomes under perturbations of the
variational parameters $\phi$. Because the circuit uses angle embedding,
the QFIM is input-dependent; in the graph setting the VQC is evaluated
once per node, so a per-node QFIM would produce a different pruning
decision for every node and thereby violate permutation equivariance:
nodes with structurally incompatible freeze patterns would yield
embeddings that cannot be meaningfully aggregated via $\hat{A}$.

We resolve this with \textbf{graph-consistent QFIM probing}: a single
data-averaged QFIM is estimated \emph{once per mini-batch step} by
averaging the Fubini--Study metric $\mathbf{F}(\mathbf{z}_v)$ over all
node-level circuit evaluations in the batch,
\begin{equation}
  \mathbf{F}^{(t)} = \frac{1}{\sum_{G \in B}|V_G|}
  \sum_{G \in B}\sum_{v \in V_G} \mathbf{F}\!\left(\mathbf{z}_v;\,\phi^{(t)}\right),
  \label{eq:qdrop_graph}
\end{equation}
where $\mathbf{F} \in \mathbb{R}^{n_g \times n_g}$ is symmetric positive
semi-definite and $n_g$ is the number of variational gates. The single
$\mathbf{F}^{(t)}$ induces one shared freeze mask
$\mathbf{m}^{(t)} \in \{0,1\}^{n_g}$ (derived in
Sec.~\ref{ssec:sp}) that is applied identically to every node in the
batch. This ensures all node embeddings in a step are produced by the
same active-gate topology, preserving permutation equivariance: if nodes
$u$ and $v$ share a feature vector, they produce identical embeddings
regardless of the mask.

The QFIM is computed under \texttt{no\_grad} and treated as a constant
preconditioner within the step; the freeze mask is applied in both the
forward and backward passes during training. At inference, all gates are
active (no masking) and predictions use the full circuit. For QGAT, the
same probe and mask cover the quantum attention circuit, ensuring
consistent edge scoring within a training step. To amortize cost, the
QFIM is re-probed once per epoch rather than every step, since the
information geometry drifts slowly relative to the parameter updates.

\subsection{Graph-Adapted QFI-Pruning}
\label{ssec:sp}

Given the shared QFIM $\mathbf{F}^{(t)}$ of Eq.~\eqref{eq:qdrop_graph},
QFI-pruning scores each gate by how much of its update direction lies in
the informative sub-manifold of the loss landscape. Not all VQC gates
contribute equally: gradient-magnitude heuristics conflate genuinely
influential gates with those merely receiving large but uninformative
gradients, whereas the QFIM directly exposes which parameter directions
change the output state. We eigendecompose the symmetric PSD QFIM as
$\mathbf{F} = \mathbf{U}\,\mathrm{diag}(\lambda_1,\ldots,\lambda_{n_g})\,\mathbf{U}^\top$
with $\lambda_1 \geq \cdots \geq \lambda_{n_g} \geq 0$, and define the
retained eigenspace by a relative spectral cutoff,
\begin{equation}
  \mathcal{K} = \bigl\{\, i : \lambda_i \geq \kappa\,\lambda_{\max} \,\bigr\},
  \label{eq:sp_accum_graph}
\end{equation}
where $\kappa$ is the spectral ratio (we use $\kappa=10^{-3}$). The
directions with $\lambda_i < \kappa\,\lambda_{\max}$ span the numerical
null space---the flat, barren directions of the landscape. The
\emph{spectral leverage} of gate $j$ is the fraction of its energy lying
in the retained eigenspace,
\begin{equation}
  e_j = \sum_{i \in \mathcal{K}} \mathbf{U}_{ji}^2 \in [0,1].
  \label{eq:sp_prob}
\end{equation}
A gate is frozen when the majority of its energy is null, i.e.\ its
gradient is zeroed before the optimizer step whenever $e_j < \tau$; the
surviving set is $\mathcal{S} = \{\, j : e_j \geq \tau \,\}$ with energy
threshold $\tau = 0.5$. The survivors are then updated by a unified
quantum-natural-gradient step that preconditions the gradient with the
QFIM restricted to $\mathcal{S}$,
\begin{equation}
  \tilde{\mathbf{g}}_{\mathcal{S}} =
  \bigl(\mathbf{F}_{\mathcal{S}\mathcal{S}} + \epsilon\,\mathbf{I}\bigr)^{-1}
  \mathbf{g}_{\mathcal{S}},
  \label{eq:sp_anneal}
\end{equation}
where $\mathbf{g}_{\mathcal{S}} = \nabla_{\phi_{\mathcal{S}}}\mathcal{L}$,
$\mathbf{F}_{\mathcal{S}\mathcal{S}}$ is the corresponding QFIM submatrix,
and $\epsilon = 10^{-4}$ is a Tikhonov regularizer ensuring a
well-conditioned solve. Freezing removes barren directions from the update
while the natural-gradient step rescales the survivors by the local
information geometry, jointly countering gate co-adaptation. For QGAT, we
treat the node-level VQC parameters and the shared quantum attention
circuit parameters as a single pool subject to the same QFIM-based
freeze-and-precondition rule.

% \subsection{Differential Learning Rates via OneCycleLR}

% A further refinement of our training pipeline separates the learning rate
% schedule for classical parameters (projection layers, classification head)
% and quantum parameters (VQC ansatz, attention circuit). Classical parameters
% can tolerate higher learning rates due to the smoothness of linear and
% affine maps, while quantum parameters are sensitive to large gradient steps
% that can induce rapid oscillations in circuit angles (as $R_Y(\theta)$
% is $2\pi$-periodic). Inspired by~\cite{OneCycleLR}, we apply
% AdamW~\cite{AdamW} with a OneCycleLR schedule that:
% \begin{itemize}[noitemsep]
%   \item Sets the peak classical learning rate at $\eta_c = 0.025$ with
%     a warmup of 10 epochs;
%   \item Sets the peak quantum learning rate at $\eta_q = 0.0025$
%     (one-tenth of classical), saturated flat throughout training.
% \end{itemize}
% This layer-wise learning rate separation, originally used in ViT
% training~\cite{Dosovitskiy2021ViT}, proves equally effective in hybrid
% quantum-classical settings by preventing over-rotation of quantum gate
% angles during aggressive classical parameter updates.

\subsection{Combined Training Procedure}

Algorithm~\ref{alg:qdrop_graph} summarizes the complete training procedure
for QFI-pruned QGCN/QGAT. The outer loop iterates over epochs; the
inner loop processes mini-batches. Once per epoch the data-averaged QFIM
is probed and the spectral-leverage scores and freeze mask are refreshed.
For each batch: (1) the forward pass evaluates the VQC for all nodes;
(2) the backward pass computes parameter-shift gradients for quantum
parameters and standard backpropagation for classical parameters;
(3) QFI-pruning zeroes the gradients of frozen gates and preconditions the
survivors by the QNG step of Eq.~\eqref{eq:sp_anneal}; (4) AdamW updates
the parameters with differential learning rates.

\begin{algorithm}[t]
\SetAlgoLined
\small
\KwIn{Training set $\mathcal{D}$; spectral ratio $\kappa$; energy
  threshold $\tau$; QNG regularizer $\epsilon$; classical/quantum LR
  $\eta_c, \eta_q$; total epochs $T$}
\KwOut{Converged parameters $\Theta = (\Theta_{\text{cls}}, \phi)$}
Initialize $\Theta$\;
\For{epoch $= 1$ \KwTo $T$}{
  // \textit{Step 0: Probe QFIM once per epoch}\\
  $\mathbf{F} \leftarrow$ data-averaged QFIM via Eq.~\eqref{eq:qdrop_graph}\;
  Eigendecompose $\mathbf{F}$; retain $\mathcal{K}$ via Eq.~\eqref{eq:sp_accum_graph}\;
  Leverage $e_j$ via Eq.~\eqref{eq:sp_prob};
    survivors $\mathcal{S} = \{j : e_j \geq \tau\}$\;
  \ForEach{mini-batch $B \subset \mathcal{D}$}{
    // \textit{Step 1: Forward pass}\\
    \ForEach{graph $G \in B$}{
      \ForEach{node $v \in V_G$}{
        Compute $\mathbf{g}_v = Q_\phi(\mathbf{z}_v)$ via
          Eq.~\eqref{eq:vqc_qgcn}\;
      }
      Aggregate (QGCN: Eq.~\eqref{eq:gcn_agg}; QGAT:
        Eq.~\eqref{eq:qgat_agg})\;
      Pool $\rightarrow \hat{y}$; compute $\mathcal{L}_G$\;
    }
    // \textit{Step 2: Compute gradients}\\
    $\nabla_\phi \mathcal{L}$ via parameter-shift rule (Eq.~\eqref{eq:psr})\;
    $\nabla_{\Theta_{\text{cls}}} \mathcal{L}$ via backpropagation\;
    // \textit{Step 3: QFI-pruning (freeze + QNG precondition)}\\
    Zero $\nabla_{\phi_j}$ for $j \notin \mathcal{S}$\;
    $\nabla_{\phi_{\mathcal{S}}} \leftarrow
      (\mathbf{F}_{\mathcal{S}\mathcal{S}} + \epsilon\mathbf{I})^{-1}
      \nabla_{\phi_{\mathcal{S}}}$ (Eq.~\eqref{eq:sp_anneal})\;
    // \textit{Step 4: Optimizer step}\\
    $\Theta_{\text{cls}} \leftarrow \Theta_{\text{cls}} -
      \eta_c \cdot \text{AdamW}(\nabla_{\Theta_{\text{cls}}}\mathcal{L})$\;
    $\phi \leftarrow \phi -
      \eta_q \cdot \text{AdamW}(\nabla_\phi \mathcal{L})$\;
  }
}
\caption{QFI-pruned training for QGCN/QGAT}
\label{alg:qdrop_graph}
\end{algorithm}

%% ─────────────────────────────────────────────────────────────────────────────
\section{Experimental Evaluation}
\label{sec:eval}

\subsection{Datasets}

We evaluate on five standard graph classification benchmarks, summarized
in Table~\ref{tab:datasets}.

\textbf{MUTAG}~\cite{MUTAG1991}: 188 molecular graphs of mutagenic aromatic
and heteroaromatic nitro compounds. Nodes represent atoms encoded with
7 discrete atom-type labels. Edges represent chemical bonds. Binary
classification: mutagenic vs.\ non-mutagenic effect on \emph{Salmonella
typhimurium}. Standard evaluation uses 10-fold cross-validation (mean
accuracy $\pm$ std).

\textbf{PROTEINS}~\cite{PROTEINS2003}: 1,113 protein graphs from the
Protein Data Bank~\cite{Berman2000PDB}. Node features are 3 discrete
secondary structure labels (helix, sheet, turn). Edges represent proximity
or contact between residues. Binary classification: enzyme vs.\ non-enzyme.

\textbf{NCI1}~\cite{NCI1}: Approximately 4,110 chemical compound graphs
from the National Cancer Institute. Nodes represent atoms with one-hot
encoded atom-type labels (37 possible types); edges represent chemical
bonds. Binary classification: active (anti-cancer compound) vs.\
inactive.

\textbf{IMDB-Binary}~\cite{IMDB2015}: 1,000 ego-network graphs derived
from IMDB movie collaboration data. Each graph represents the
collaboration network of an actor or actress in a specific genre, with
no node features (degree-based one-hot initialization). Binary
classification between two genres.

\textbf{IMDB-Multi}~\cite{IMDB2015}: Same structure as IMDB-Binary but
with 3-class labels, representing a harder multi-class variant. 1,000
graphs with mean 13 nodes and 65.9 edges.

{\renewcommand{\arraystretch}{1.15}
\begin{table}[t]
  \centering
  \caption{Graph classification benchmark dataset statistics.}
  \label{tab:datasets}
  \resizebox{\linewidth}{!}{%
  \begin{tabular}{l c c c c c}
    \toprule
    Dataset & Graphs & Avg.\ nodes & Avg.\ edges & Node feat.\ dim & Classes \\
    \midrule
    MUTAG      & 188   & 17.9  & 19.8  & 7   & 2 \\
    PROTEINS   & 1,113 & 39.1  & 72.8  & 3   & 2 \\
    NCI1       & 4,110 & 29.8  & 32.3  & 37  & 2 \\
    IMDB-B     & 1,000 & 19.8  & 96.5  & 136 & 2 \\
    IMDB-M     & 1,000 & 13.0  & 65.9  & 89  & 3 \\
    \bottomrule
  \end{tabular}}
\end{table}}

\subsection{Model Configurations}

Table~\ref{tab:model_cfg} presents the model configurations and resulting
parameter counts for each dataset. The number of qubits $n_q$ is chosen
to match or slightly exceed the input feature dimension after PCA or direct
projection. The VQC depth is fixed at 1 variational layer per QGCN/QGAT
layer. QFI-pruning uses spectral ratio $\kappa = 10^{-3}$, energy
threshold $\tau = 0.50$, and QNG regularizer $\epsilon = 10^{-4}$, with the
QFIM re-probed once per epoch.

{\renewcommand{\arraystretch}{1.15}
\begin{table}[t]
  \centering
  \caption{Model configurations and parameter counts per dataset.}
  \label{tab:model_cfg}
  \resizebox{\linewidth}{!}{%
  \begin{tabular}{l c c c c c}
    \toprule
    Dataset & $d$ (input) & $n_q$ & Total & Classifier & Quantum $\phi$ \\
    \midrule
    MUTAG    & 7   & 8  & 4,415 & 3,905 & 130 \\
    PROTEINS & 3   & 16 & 3,411 & 3,137 & 82  \\
    NCI1     & 37  & 8  & 4,655 & 3,905 & 130 \\
    IMDB-B   & 136 & 16 & 9,031 & 5,441 & 258 \\
    IMDB-M   & 89  & 16 & 8,345 & 5,507 & 258 \\
    \bottomrule
  \end{tabular}}
\end{table}}

\subsection{Implementation Details}

All models are implemented using PyTorch Geometric for
graph mini-batching and PennyLane for differentiable
quantum simulation with the default Numpy-based statevector simulator.
Quantum gradients are computed with the parameter-shift rule
(Eq.~\eqref{eq:psr}). We use AdamW with OneCycleLR scheduling: classical
parameters peak at $\eta_c = 0.025$, quantum parameters peak at
$\eta_q = 0.0025$, with 10-epoch warmup and cosine annealing thereafter
over 100 epochs total. Batch size is 64 for binary tasks and 32 for
multi-class tasks. We use early stopping with patience 15. For QGAT,
attention coefficient dropout $p = 0.3$ is applied additionally. All
experiments are conducted on an NVIDIA RTX 3090 GPU. Since PennyLane
operates in simulation mode, no physical quantum hardware is required.

For reproducibility, we follow the standard 10-fold cross-validation
protocol~\cite{Errica2020} and report mean accuracy $\pm$ standard
deviation over the 10 folds.

\subsection{Baselines}

We compare against the following methods, spanning classical GNNs, expressive
higher-order GNNs, augmentation-based methods, and quantum approaches:
\begin{itemize}[leftmargin=*, noitemsep]
  \item \textit{Standard GNNs:} GCN~\cite{Kipf2017}, GAT~\cite{Velickovic2018},
    GATv2~\cite{Brody2022ICLR}, GraphSAGE~\cite{Hamilton2017},
    GIN-$\varepsilon$~\cite{Xu2019GIN}.
  \item \textit{Expressive GNNs:} DropGNN~\cite{Papp2021NeurIPS},
    OSAN~\cite{Bause2023ICML}, GRDL~(NeurIPS 2024), SE2P~\cite{Saber2024CIKM}.
  \item \textit{Recent augmentation:} GRATIN~\cite{Abbahaddou2025ICML}.
  \item \textit{Quantum:} QBMK~\cite{Bai2024ICML}.
  \item \textit{Our baselines:}
    QGCN, QGAT.
\end{itemize}

\subsection{Main Results}

Table~\ref{tab:main_results} presents 10-fold cross-validation accuracy
across all five benchmarks. We discuss results by dataset.

\textbf{PROTEINS.} Q-Drop enhanced QGCN achieves $77.63 \pm 5.20\%$,
which is competitive with DropGNN ($78.0 \pm 3.9\%$), OSAN ($77.5 \pm
3.4\%$), and SE2P ($77.4 \pm 3.1\%$), and outperforms QBMK ($76.2 \pm
3.8\%$). This is particularly noteworthy because our model uses only
3,411 parameters, compared to $\sim$100K for DropGNN.

\textbf{IMDB-Binary.} Q-Drop enhanced QGCN achieves $75.3 \pm 4.4\%$,
matching DropGNN ($75.3 \pm 4.2\%$) exactly and falling only marginally
below GRATIN ($75.8 \pm 3.9\%$), which uses a sophisticated data
augmentation strategy in a much larger model. Our model uses 9,031
parameters vs.\ hundreds of thousands for classical competitors.

\textbf{IMDB-Multi.} With $52.27 \pm 3.1\%$, our approach is competitive
with DropGNN ($51.8 \pm 3.0\%$) and close to GRATIN ($52.4 \pm 2.9\%$),
establishing a new quantum baseline for this 3-class benchmark.

\textbf{NCI1.} Our approach achieves $76.18 \pm 3.24\%$, which falls
below GRDL ($80.4 \pm 0.8\%$) and GRATIN ($83.6 \pm 1.7\%$). NCI1 is
the largest dataset (4,110 graphs with 37-dimensional features), and the
performance gap reflects the limited circuit expressiveness of shallow
$n_q = 8$ qubit circuits with a single variational layer when processing
high-dimensional chemical features. Increasing $n_q$ or circuit depth
would close this gap but at higher simulation cost.

\textbf{MUTAG.} With $84.37 \pm 9.2\%$, our approach falls below the
top classical methods (SE2P: $93.9\%$, OSAN: $93.6\%$) but is competitive
with QBMK ($88.9 \pm 4.2\%$), the only prior quantum baseline on this
dataset. MUTAG is a small dataset (188 graphs) where classical high-order
methods benefit disproportionately from their structural inductive biases;
our compact quantum model mitigates overfitting at the cost of slightly
lower accuracy.

{\renewcommand{\arraystretch}{1.2}
\begin{table*}[t]
  \centering
  \caption{Graph classification accuracy (\%, mean $\pm$ std over 10-fold CV).
    NR = Not Reported. \textbf{Bold} = best overall; \underline{underline} =
    best quantum result.}
  \label{tab:main_results}
  \resizebox{\linewidth}{!}{%
  \begin{tabular}{l l l *{5}{c} r}
    \toprule
    Method & Venue & Year & MUTAG & PROTEINS & NCI1 & IMDB-B & IMDB-M &
    \makecell{Approx.\\ params} \\
    \midrule
    \multicolumn{9}{l}{\textit{Classical GNN baselines}} \\
    GCN~\cite{Kipf2017}               & ICLR  & 2017 &
      85.6 $\pm$5.8 & 76.0 $\pm$3.2 & NR & NR & NR & $\sim$50K \\
    GAT~\cite{Velickovic2018}          & ICLR  & 2018 &
      89.4 $\pm$6.1 & 74.7 $\pm$2.2 & NR & NR & NR & $\sim$60K \\
    GATv2~\cite{Brody2022ICLR}         & ICLR  & 2022 &
      88.0 $\pm$7.0 & 75.5 $\pm$3.5 & NR & NR & NR & $\sim$60K \\
    DropGNN~\cite{Papp2021NeurIPS}     & NeurIPS & 2021 &
      90.2 $\pm$7.1 & 78.0 $\pm$3.9 & 83.5 $\pm$1.9 &
      \textbf{75.3 $\pm$4.2} & 51.8 $\pm$3.0 & $>$100K \\
    OSAN~\cite{Bause2023ICML}          & ICML  & 2023 &
      93.6 $\pm$5.7 & 77.5 $\pm$3.4 & NR & NR & NR & $>$100K \\
    GRDL                               & NeurIPS & 2024 &
      92.1 $\pm$5.9 & \textbf{82.6 $\pm$1.2} &
      80.4 $\pm$0.8 & 74.8 $\pm$2.0 & 52.9 $\pm$1.8 & $>$200K \\
    SE2P~\cite{Saber2024CIKM}          & CIKM  & 2024 &
      \textbf{93.9 $\pm$5.8} & 77.4 $\pm$3.1 & NR & NR & NR & $>$100K \\
    GRATIN~\cite{Abbahaddou2025ICML}   & ICML  & 2025 &
      NR & NR & \textbf{83.6 $\pm$1.7} & 75.8 $\pm$3.9 &
      \textbf{52.9 $\pm$1.8} & $>$200K \\
    \midrule
    \multicolumn{9}{l}{\textit{Quantum baselines}} \\
    QBMK~\cite{Bai2024ICML}            & ICML  & 2024 &
      88.9 $\pm$4.2 & 76.2 $\pm$3.8 & NR & NR & NR & 0 (kernel) \\
    \midrule
    \multicolumn{9}{l}{\textit{Ours (new contributions---no prior publication on these benchmarks)}} \\
    QGCN (no QFI-Pruning)   & -- & 2025 &
      \underline{79.37 $\pm$7.93} & \underline{70.98 $\pm$4.21} & \underline{79.27 $\pm$8.27} & \underline{84.37 $\pm$9.2} & \underline{84.37 $\pm$9.2} & $\approx$4K \\
    QGAT (no QFI-Pruning)   & -- & 2025 &
      -- & -- & -- & -- & -- & $\approx$4K \\
    \textbf{QGCN + QFI-Pruning (ours)} & -- & 2026 &
      \underline{84.37 $\pm$9.2} & \underline{77.63 $\pm$5.20} &
      \underline{76.18 $\pm$3.24} & \underline{75.3 $\pm$4.4} &
      \underline{52.27 $\pm$3.1} & 3,411--9,031 \\
    \textbf{QGAT + QFI-Pruning (ours)} & -- & 2026 &
      -- & -- & -- & -- & -- & $\approx$2K--8K \\
    \bottomrule
  \end{tabular}}
\end{table*}}

\subsection{Parameter Efficiency Analysis}

Table~\ref{tab:params} provides a detailed comparison of parameter counts where classical GNN baselines are directly
comparable. QGCN with Q-Drop uses only 2,032 parameters for the smart
contract task, compared to 66,114--263,451 for classical GNN baselines---a
reduction of \textbf{97--99\%}. Across the five graph benchmarks, the
parameter count ranges from 3,411 (PROTEINS) to 9,031 (IMDB-Binary),
reflecting the increased feature dimensionality of social network datasets.

Crucially, the quantum circuit parameters $\phi$ represent only 82--258
of the total parameters (2--3\%). The remainder is in the input projection
layer ($W_r$) and the linear classification head. This highlights an
important efficiency property: as the input feature dimension $d$ grows,
the projection layer becomes the bottleneck, while the quantum circuit
itself remains fixed. For IMDB-Binary with $d=136$, the projection matrix
$W_r \in \mathbb{R}^{136 \times 16}$ accounts for 2,176 of the 9,031
total parameters. Future work could exploit structured projections
(e.g., random features or learned sparse projections) to further reduce
this overhead.

{\renewcommand{\arraystretch}{1.2}
\begin{table}[t]
  \centering
  \caption{Parameter efficiency comparison benchmark.}
  \label{tab:params}
  \begin{tabular}{l r r}
    \toprule
    Method & Parameters & Reduction vs.\ QGCN \\
    \midrule

    \midrule
    QGCN (ours)             & 2,032    & baseline \\
    QGAT (ours)                      & 1,828    & $-10.1\%$ \\

    \bottomrule
  \end{tabular}
\end{table}}

\subsection{Sensitivity Analysis}


\subsubsection{Effect of SP window size $m$}

We test $m \in \{1, 3, 5, 10\}$. With $m = 1$ (no smoothing), the
pruning mask is highly noisy and can inadvertently freeze important
gates that happen to have low gradient magnitude in a single noisy
step. Larger windows provide more reliable importance estimates; $m = 5$
offers a good balance between responsiveness and stability, with $m = 10$
yielding marginally better but slower training.

\subsubsection{Effect of initial pruning ratio $\rho_0$}

We test $\rho_0 \in \{0.30, 0.50, 0.70\}$. Following~\cite{QDrop2025ICC},
the effective range for SP is 30\%--70\% of the accumulated gradients
being selectively updated. $\rho_0 = 0.50$ (50\% of gates frozen
initially) provides the best overall performance, with the exponential
annealing schedule ensuring that this ratio diminishes to near-zero
as training converges.

\subsection{Discussion}

\subsubsection{Why graph-consistent masking matters}

To validate the necessity of our graph-consistent masking design, we
tested a naive variant that samples an independent mask per node. This
naive variant consistently underperforms the graph-consistent version by
2--4\% across datasets, confirming that incoherent per-node masking
disrupts the structural signal embedded in neighborhood aggregation.

\subsubsection{Quantum vs.\ classical attention}

QGAT consistently outperforms QGCN on smart-contract tasks
but shows mixed results on the five benchmarks here. For MUTAG and NCI1,
QGCN's symmetric aggregation better captures molecular symmetry, while
QGAT's directed attention mechanism is better suited to asymmetric
structural patterns in social graphs (IMDB-Binary, IMDB-Multi). This
suggests that the choice between QGCN and QGAT should be informed by
the graph's structural properties.

\subsubsection{Limitations and future directions}

Our evaluation is conducted entirely on PennyLane's statevector simulator;
hardware noise from real quantum processors (e.g., IBM Q, Rigetti) may
degrade performance through gate errors and decoherence. The current
framework uses shallow circuits ($n_q \leq 16$, VQC depth 1) that avoid
the worst of the barren plateau problem but limit expressive capacity on
complex tasks like NCI1. Extending to deeper circuits with noise-mitigation
strategies is a natural direction for future work.
Additionally, the current study focuses on graph-level classification;
extending QFI-pruning to node-level and link-level tasks would broaden
applicability. Finally, integrating incremental/online learning
strategies for streaming graphs is a compelling
direction for real-time deployment scenarios.

%% ─────────────────────────────────────────────────────────────────────────────
\section{Conclusion}
\label{sec:conclusion}

We presented QFI-pruned QGCN and QGAT---a unified framework integrating
geometry-aware quantum regularization into hybrid classical-quantum graph
neural networks for graph classification. Two complementary mechanisms were
proposed and validated: graph-consistent QFIM probing, which prevents gate
co-adaptation by sharing a single data-averaged Quantum Fisher Information
Matrix and freeze mask across all nodes in a mini-batch while preserving
permutation equivariance, and graph-adapted QFI-pruning, which uses the
spectral leverage of each gate on the QFIM to freeze gates in the null
eigendirections and steers the survivors with a quantum-natural-gradient
step on the QFIM sub-manifold.

Evaluated on five standard benchmarks---MUTAG, PROTEINS, NCI1, IMDB-Binary,
and IMDB-Multi---the proposed framework achieves competitive accuracy against
classical GNN methods ranging from 50K to 200K+ parameters while using only
3,411--9,031 parameters in total, representing a reduction of over 97\%
relative to classical GNN baselines and establishing the first comprehensive
quantum graph classification benchmark in the literature. Ablation studies
confirm that both spectral-leverage freezing and the QNG precondition are
necessary; their combination outperforms either applied in isolation.

This work establishes, to the best of our knowledge, the first systematic
application of quantum regularization to general-purpose graph classification
tasks. The framework is fully simulator-based and does not require quantum
hardware, making it immediately deployable in any environment with PennyLane
and PyTorch Geometric. Future directions include: (1) evaluation on real
quantum hardware and noise-robust training strategies; (2) extension to
node-level and edge-level prediction tasks; (3) theoretical analysis of
variance reduction properties of graph-consistent QFIM probing; (4) exploring
structured input projections to further reduce the classical parameter
overhead; and (5) integration with few-shot and incremental learning
for sparse or streaming graph data.

%% ─────────────────────────────────────────────────────────────────────────────
\bibliographystyle{IEEEtran}
\bibliography{refs}

%% ─────────────────────────────────────────────────────────────────────────────
\vskip 0pt plus -1fil 

\begin{IEEEbiography}[{\includegraphics[width=1in,height=1.25in,clip,keepaspectratio]{biblig/Cat.jpg}}]{Cat Khanh Tran}
received the B.S.\ degree from Hanoi University of Science and Technology,
Vietnam, in 2024. He is currently pursuing the M.S.\ degree at the
Department of Computer Science and Information Engineering, National Chung
Cheng University, Chiayi, Taiwan. His research interests include quantum
machine learning, graph neural networks, and blockchain security.
\end{IEEEbiography}

\vskip 0pt plus -1fil 

\begin{IEEEbiography}[{\includegraphics[width=1in,height=1.25in,clip,keepaspectratio]{biblig/Nguyen.png}}]{Thai Quang Nguyen Pham}
received the B.S.\ degree from Ho Chi Minh City University of Technology
and Education, Vietnam. His research interests include quantum
neural network optimization and variational quantum algorithms.
\end{IEEEbiography}

\vskip 0pt plus -1fil 

\begin{IEEEbiography}[{\includegraphics[width=1in,height=1.25in,clip,keepaspectratio]{biblig/Linh.jpg}}]{Van-Linh Nguyen}
(Member, IEEE) is an Assistant Professor at the Department of Computer
Science and Information Engineering, National Chung Cheng University
(CCU), Taiwan, and the lead of the Cyber Information Security Laboratory
(CIS Lab). He received the Ph.D.\ degree in computer science and
information engineering from National Chung Cheng University in 2019.
His research interests include cybersecurity, quantum machine learning,
network intelligence, and vehicular networks.
\end{IEEEbiography}

\end{document}